\begin{abstract}

This report investigates the dynamic behaviour of a car's suspension system, and the effect that varying the spring and dampener coefficients has on the system's ability to absorb forces from traversing a bumpy road.

The system was analysed by first simplifying it, by modelling the car as a beam, and the wheels as points with a spring and a dampener each connected to the beam. The passengers were modeled by pendulums with torsion springs keeping them upright. Then Newtonian Mechanics were used to derive a coupled system of differential equations, and then solved numerically using the Runge-Kutta method.

The effects of varying the spring coefficient $k$ and the dampening coefficient $c$ were analysed. The system exhibited minimal vibrations for specific $k$ and $c$ values. This demonstrates that there are optimal parameters to make vehicles comfortable and smooth to drive, and modelling the system like this with specific parameters can be used as a starting point for suspension design.
\end{abstract}